%-----------------------------------------------------------------------
\section{The current way}
\label{sec:current}
%-----------------------------------------------------------------------

The comparison in this paper is against three solvers in current practice:
PBRT-v4's stochastic progressive photon mapping, Blender Cycles' path tracer,
and Mitsuba~3. They were chosen because they are the tools a practitioner would
reach for, because two of them are open and inspectable, and because between
them they cover both estimator families of Section~\ref{sec:cg}. This section
states what each was asked to do, what it cost, and what it produced.

\subsection{The solvers}
\label{sec:solvers}

\paragraph{PBRT-v4 SPPM.} PBRT \cite{pharr2023} is the reference implementation
of the book's rendering algorithms and its SPPM integrator
\cite{hachisuka2008} is the strongest conforming caustic solver available
without special-purpose machinery. It is the solver we treat as the
\emph{reference conforming method}: it estimates the same integrand as we do,
it does not use wave optics, and it is consistent in the limit. We configure
SPPM with 100 iterations, $5 \times 10^5$ photons per iteration ($5 \times 10^7$
photons total), an initial search radius of $R_0 = \SI{5}{\milli\meter}$, and
radius reduction exponent $\alpha = 0.67$. We generate the
scene programmatically from the same geometry the other renderers receive, so
that differences between the two are attributable to the estimator rather than
to the scene.

\paragraph{Blender Cycles.} Cycles is a production path tracer and the renderer
used by the showcase assets of this project. It represents the backward family
of Section~\ref{sec:backward}. Its caustic handling is enabled through the
standard path-guiding and shadow caustics options, rendered at 2048 samples
per pixel with up to 12 light bounces. The isolate used for scoring is formed as
caustics-on minus caustics-off, which is the positive-part differential defined
in Section~\ref{sec:estimators}.

\paragraph{Mitsuba 3.} Mitsuba \cite{jakob2022} is a research renderer with a
spectral mode and well-documented integrators, and it was used as a third
opinion on the transport physics rather than as a primary caustic competitor.
It was configured with the spectral path tracer (`ptracer`, 32 spectral bins
from \SI{360}{\nano\meter} to \SI{700}{\nano\meter}, 1024 samples per pixel).
We report it where relevant and do not put it in the quality table, because its
spectral sampling makes its noise behaviour differ from the three-channel
solvers in a way that would confound a like-for-like comparison.

\subsection{One scene, one source, one set of materials}
\label{sec:protocol}

A comparison of caustic renderers is easy to make meaningless. If the geometry,
the materials or the source differ between two methods, the numbers describe
the difference between scenes rather than between estimators. Three things have
to be identical or the comparison does not mean anything.

\begin{enumerate}
\item \textbf{Geometry.} The same water-surface vertices and the same cylinder
      are exported to every renderer by one exporter, so the body is identical
      and the caustic is at the same place in all of them.
\item \textbf{Materials.} Every renderer receives the measured indices of
      Section~\ref{sec:dispersion}. Where a renderer cannot carry dispersion in
      the transport itself, it is given the central-channel index and this is
      stated; assigning a different index to one side would move its caustic.
\item \textbf{Source and receiver.} One distant source of the stated angular
      size, one matte receiver patch of the stated size, no environment light
      added to one side and not another for the scored comparisons.
\end{enumerate}

\subsection{What the current way costs}
\label{sec:current-cost}

Table~\ref{tab:cost} reports the budget each solver needed to produce a usable
caustic on this scene, against our bake. The comparison is between the cost of
the caustic and the cost of the frame, because that is what a user pays; quoting
a photon-mapping pass alone would understate the reference solvers.

\begin{table}[tbp]
\centering
\caption{Cost to produce the caustic. Our CPU figures are measured on the host
workstation (8-core CPU); the reference solvers are measured on the same machine.
The SPPM figure is the accumulation time needed for the fold to emerge above noise.
The wave correction row reports the wave term alone; geometric quadrature plus wave
correction equals the full single-thread bake ($28.7 + 49.8 = \SI{78.5}{\milli\second}$).
Section~\ref{sec:ours-cost} explains why the comparison is stated per frame
rather than per photon.}
\label{tab:cost}
\small
\begin{tabular}{@{}lrrr@{}}
\toprule
Quantity & Ours & SPPM & Cycles \\
\midrule
Grid / resolution                  & $512^2$ & adaptive & $512^2$ \\
Geometric quadrature (1 thread, CPU) & \SI{28.7}{\milli\second} & \SI{34}{\second} & --- \\
Wave correction pass (1 thread, CPU) & \SI{49.8}{\milli\second} & n/a & n/a \\
Full bake (1 thread, CPU)          & \SI{78.5}{\milli\second} & --- & --- \\
Full bake (8 threads, CPU)         & \SI{36.2}{\milli\second} & --- & --- \\
8-thread speedup                   & $2.17\times$ & --- & --- \\
Iterative / noisy                  & no & yes & yes \\
\bottomrule
\end{tabular}
\end{table}

The cost structure of our own method is worth stating because it is the obvious
target for anyone continuing this work: the Airy correction is roughly twice the
quadrature pass, so two thirds of a bake is the wave term. That term is
embarrassingly parallel per channel and is currently dominated by the exact
Euclidean distance transform used to get the signed distance to the band.

\subsection{What the current way produces}
\label{sec:current-quality}

Three measured properties of the reference solvers motivate the method in
Section~\ref{sec:ourway}. Each was measured on the same scene and with the
conventions of Section~\ref{sec:estimators}.

\paragraph{Kernel smoothing versus physical diffraction.} The photon-map fold
width and brightness track the kernel radius $h$ rather than the physics, as
analyzed in Section~\ref{sec:photonmap} and Lemma~\ref{lem:kdebias}. SPPM reaches a finer feature scale than
a fixed-radius photon map by progressively shrinking $h$, but pays for it in
iteration count: the \SI{34}{\second} of Table~\ref{tab:cost} is the cost of
reducing the kernel footprint sufficiently for the ridge to emerge from background noise.
In contrast, our wave-corrected fold replaces geometric divergence with the
physically bounded Airy scale $w \propto \lambda^{2/3}$ without kernel blur.

\paragraph{Coverage of the caustic core.} Measured against the independent
ground truth's core region, Cycles' caustic covers \SI{54.7}{\percent} of it
where ours covers \SI{99.9}{\percent}. This is a statement about how much of the
feature a practical finite-sample frame actually reconstructs, where the path
tracer's deficit is concentrated near the fold where backward paths are least
likely to sample the singular specular manifold.

\paragraph{Temporal instability.} For an animation, the primary visual defect
of a stochastic caustic is flicker rather than stationary grain, and
Section~\ref{sec:results-flicker} measures it. A deterministic forward bake
naturally avoids Monte Carlo sample chattering; our caustic exhibits
\SI{0.001}{} levels of mean frame-to-frame change against \SI{0.010}{} for the
path-traced clip, a factor of \num{7.9}. Both clips were denoised, which
suppresses flicker in both and therefore confirms the stability of forward
quadrature under camera and surface motion.

SPPM and Cycles are general renderers solving a harder problem than the one
posed here, and both are consistent within their model. The numbers above say
that the geometric model cannot represent a fold cheaply, which is a statement
about the model. A reader who needs that model's other features should use those
renderers.

\subsection{Shape agreement with the reference solver}
\label{sec:current-shape}

Independent of the ground-truth axis, the converged two-dimensional shape of our
field and the SPPM field agree at $r = +0.932$ on 40\,mm tiles. This is a
necessary check rather than an impressive one: two methods that disagree on
shape cannot both be describing the same caustic, and a high correlation here
combined with a lower one against the ground truth would have indicated a shared
error. Since SPPM also scores $0.853$ on the independent axis and ours scores
$0.917$, the two agree with each other and both agree with the trace, with ours
closer.

\subsection{Mitsuba as a third opinion}
\label{sec:current-mitsuba}

Mitsuba was run in spectral mode on the same scene. Its role in this paper is to
confirm the transport and the dispersive behaviour rather than to compete on
caustic quality: its spectral sampling spreads energy across wavelength bins,
so its per-channel noise is not comparable to a three-channel renderer's, and
scoring it in the quality table would compare noise models rather than
estimators. We note it because a reader is entitled to ask whether the physics
of the scene was checked by an implementation with a different architecture, and
this is where that check lives.
